\section*{3 Main results: Inner-product consequences}
Our main technical result, proved in Section5, states that a ``large'' collection of functions that are ``close'' to being orthonormal, is computationally hard for SM protocols in which the referee is an unbounded fan-in circuit of constant depth, and one player is limited to ``short'' preprocessing output length.
One may wonder about the necessity of satisfying the one-to-one condition. While it may still be the case that the same result (or even better) follows without this assumption, if we could obtain it that way, then we could easily exhibit a function with stronger lower bounds. As an example, consider taking the inner product function and computing it only on a prefix of the input while ignoring other bits--relying on the main theorem's consequence, we could extend it to an arbitrary $\Omega(n)$ lower bound on the preprocessing output length.
The (somewhat cumbersome) second requirement on a function $f:\{-1,1\}^n\times\{-1,1\}^n\to\{-1,1\}$ specified in Theorem9 can be replaced by a slightly stronger requirement, yet one involving the more familiar measure borrowed from the study of learning by statistical queries~\cite{ref23}. Define the statistical query dimension of $f$ with respect to the uniform distribution to be the size of the largest set $D\subseteq\{-1,1\}^n$ such that $|\langle f_x,f_{x'}\rangle|\leq1/|D|$ for every $x\ne x'\in D$. More details on statistical query dimension can be found in~\cite{ref7}.
A simple consequence of Theorem9 is Theorem1. Here we state a more general version that refers to statistical query dimension. At a high level, the theorem says that computing with preprocessing a function having exponential statistical query dimension remains as hard for $\operatorname{AC}^0$ as without, given that one player is limited to output a string whose length stretches the input length by an additive sublinear term:
\setcounter{nttheorem}{9}
\begin{ntproposition}[Author Proposition10: Formal version of Theorem1]
Let $k=o(n)$, and suppose that a function $f:\{-1,1\}^n\times\{-1,1\}^n\to\{-1,1\}$ satisfies the one-to-one condition and has statistical query dimension of $2^{\Omega(n)}$. If $f$ admits an SM protocol $(A,B,C)$ such that $B:\{0,1\}^n\to\{0,1\}^{n+k}$ and $C$ is an unbounded fan-in circuit of depth $h$ and size $M$, then:
\[M\geq2^{\Omega_h([n/k]^{1/(h-1)})}.\]
\end{ntproposition}
\begin{proof}
Let $D\subseteq\{-1,1\}^n$ be a set of size $2^{\Omega(n)}$ such that $|\langle f_x,f_{x'}\rangle|\leq1/|D|$ for every $x\ne x'\in D$. One can easily find an $0<\alpha\leq1/2$ for which $H(\alpha)$ is small enough, such that setting $t=\alpha(n+k)$ gives
\[13\cdot2^{2(k+1)}\cdot\binom{n+k}{\leq\alpha(n+k)}\ \underset{\text{Lemma8}}\leq\ 13\cdot2^{2(k+1)}\cdot2^{H(\alpha)(n+k)}\leq|D|/6.\]
Now, let $X\subseteq D$ of size $|X|=13\cdot2^{2(k+1)}\cdot\binom{n+k}{\leq t}$. We have:
\begin{align*}\mathbb E_{x\ne x'\sim X}[\langle f_x,f_{x'}\rangle^2]&\leq\mathbb E_{x\ne x'\sim D}[\langle f_x,f_{x'}\rangle^2]\leq\left(\mathbb E_{x\ne x'\sim D}[\langle f_x,f_{x'}\rangle]\right)^2\\&\leq\frac1{|D|^2}\leq\frac1{36|X|^2}\leq\frac{2^{2k}}{36|X|^2}.\end{align*}
Thus, by Theorem9,
\[M\geq2^{\Omega_h([t/k]^{1/(h-1)})}=2^{\Omega_h([\alpha(n+k)/k]^{1/(h-1)})}=2^{\Omega_h([n/k]^{1/(h-1)})}.\]
\end{proof}
Since $\operatorname{IP}$ satisfies the one-to-one condition and has the largest possible statistical query dimension of $2^n$, we get the following corollary.
\begin{ntcorollary}[Author Corollary11]
Let $k\leq n^\alpha$ for some $0\leq\alpha<1$, and suppose that $\operatorname{IP}$ admits an SM protocol $(A,B,C)$ such that $B:\{0,1\}^n\to\{0,1\}^{n+k}$ and $C$ is an unbounded fan-in circuit of depth $h$ and size $M$. Then:
\[M\geq2^{\Omega_h(n^{(1-\alpha)/(h-1)})}.\]
\end{ntcorollary}
\begin{ntcorollary}[Author Corollary12]
Let $k\leq n/\log^\beta n$ for every $\beta>0$ (for large enough $n$), and suppose that $\operatorname{IP}$ admits an SM protocol $(A,B,C)$ such that $B:\{0,1\}^n\to\{0,1\}^{n+k}$ and $C$ is an unbounded fan-in circuit of depth $h$ and size $M$. Then:
\[M\geq2^{\Omega_h(\log^c n)}\qquad\text{for every }c>0.\]
\end{ntcorollary}
